\documentclass[10pt,a4paper]{article}
\usepackage[T1]{fontenc}
\usepackage{lmodern,amsmath,amssymb,microtype,xcolor,geometry,hyperref}
\geometry{margin=19mm}
\definecolor{softgray}{gray}{0.94}
\hypersetup{hidelinks,pdftitle={SGA 2 Expose VIII proof comment and notation},pdfauthor={Interlanguage English/Germanic lane}}
\begin{document}
\begin{center}
  {\Large\bfseries SGA 2: Local Cohomology and Lefschetz Theorems}\par\smallskip
  {\large Expos\'e VIII: The Finiteness Theorem}\par
  \small Bounded source-aligned working unit --- 19 July 2026
\end{center}
\noindent\colorbox{softgray}{\parbox{0.95\linewidth}{\footnotesize\textbf{Authority.}
Corrected French lines 2583--2595; original printed pp.~86--87; physical
source-PDF p.~77; recomposed running p.~69. This bounded unit has been
independently source-reviewed against the corrected French TeX and direct scan;
cumulative integration remains pending. Line 2596 is blank; the construction
of the double complex begins at line 2597 and is excluded.}}

The proof is left to the reader.\footnote{See also H.~Cartan and
S.~Eilenberg, \emph{Homological Algebra}, Princeton Mathematical Series,
vol.~19, Princeton University Press, 1956.}

Let us recall a notation introduced in Expos\'e V.

\medskip
\noindent\textbf{Notation.} --- Let $X^\bullet$ and $Y^\bullet$ be two
complexes. We denote by $\operatorname{Hom}^\bullet(X^\bullet,Y^\bullet)$ the
\emph{simple complex} whose component of degree $n$ is
\[
  (\operatorname{Hom}^\bullet(X^\bullet,Y^\bullet))^n
  =\prod_{-p+q=n}\operatorname{Hom}(X^p,Y^q),
\]
also denoted by $\operatorname{Hom}^n(X^\bullet,Y^\bullet)$, and whose
differential is given by
\begin{gather*}
  \partial_n\colon
  \operatorname{Hom}^n(X^\bullet,Y^\bullet)
  \longrightarrow
  \operatorname{Hom}^{n+1}(X^\bullet,Y^\bullet),\\
  \partial_n=d'+(-1)^{n+1}d'',
\end{gather*}
where $d'$ and $d''$ are the differentials (of degree $+1$) induced by those
of $X^\bullet$ and $Y^\bullet$.
\end{document}
